

\illustration{0.3}{library.jpg}{
    Pixabay Content License by \\ Andrew Tan on \href{https://pixabay.com/photos/library-la-trobe-study-students-1400313/}{Pixabay}
}

% optionally define variables/limits for this problem




You have \emph{a lot} of exams today! And you have not yet prepared for any of them!
At least you know from experience that if you study enough for an exam, you
will for sure pass it quickly, well within the allotted time. Even better, you stared so long
at the daunting course curricula that you now know exactly how much time you will
need to study for each exam in order to pass it within the given time. If
you do not study long enough, you will for sure fail it.

As it happens, your university has some weird bureaucratic rules that require you
to attend \emph{all} your exams. The horror! And leaving early is not allowed,
unless you know for sure you passed it!

Given the full exam schedule of when each exam starts, how long each exam
takes, how much time you need to study for each exam, and how quickly you can
finish each exam you studied for, how many exams can you pass at most?

You may start studying at time $0$, but can only study while not making an exam.
The preparation for an exam does not have to be done in a contiguous block of
time and may be interrupted.

